\begin{korollar}{Erwartungswert des Infimums integrierbarer Funktionen}
                {EndlicheDarstellungDesInfimumsIntegrierbarerFunktionen}
Seien $I$ eine nichtleere, abzählbare Menge und 
$\folge X i I \in L_1(\Omega, \mfa,\Prim,\R)^I$.\\
Es gelte:
\begin{enumerate}[(a)]
 \item für alle $\omega \in \Omega$ ist $\menge{X_i(\omega)}{i \in I}$
       nach unten beschränkt,
 \item $\Inf{i\in I} X_i \in L_1(\Omega, \mfa,\Prim,\R)$ (wohldefiniert
       wegen (a) und $I \neq \emptyset$).
\end{enumerate}
Dann gilt:
\begin{enumerate}[(1)]
 \item\label{EndlicheDarstellungDesInfimumsIntegrierbarerFunktionen1}
      $\menge{\E\left(\Inf{i\in E}X_i\right)}{E\in\Pot_{ne}(I)}$ ist nichtleer
      und nach unten beschränkt.
 \item\label{EndlicheDarstellungDesInfimumsIntegrierbarerFunktionen2}
      $\inf\menge{\E\left( \Inf{i\in E} X_i\right)}{E\in\Pot_{ne}(I)} 
      = \E\left(\Inf{i\in I}X_i \right)$.
\end{enumerate}
\end{korollar}
\begin{proof}
Für alle $i \in I$ sei $Y_i \coloneqq  -X_i$.\\
Für alle $i \in I$ ist dann $Y_i$ als Vielfaches einer integrierbaren Funktion
auch integrierbar.\\
Sei $\omega \in \Omega$. 
Dann ist
$ \menge{Y_i(\omega)}{i \in I} = \menge{-X_i(\omega)}{i \in I}
 = -\menge{X_i(\omega)}{i \in I}$.
Daher ist $\menge{Y_i(\omega)}{i \in I}$ wegen $(a)$ nach oben beschränkt.\\
Sei $A\in\Pot_{ne}(I)$. Wegen $A \neq \emptyset$, gibt es ein $a \in A$, sodass
gilt:
\[X_a \ge \inf_{i \in A}X_i \ge \inf_{i \in I}X_i\]
und wegen (b) bedeutet das nach \refclaim{Integrierbarkeitskriterien}
{Integrierbarkeitskriterien2}
\[ \inf_{i \in A}X_i \text{ ist integrierbar, also auch }
  -\inf_{i \in A}X_i = \sup_{i \in A}Y_i\text.\]
Somit gilt:
$\forall\,E\in\Pot_{ne}(I)\cup\meng{I}:\Sup{i \in A}Y_i = -\Inf{i\in A}X_i
\text{ integrierbar}$. \hfill $(\ast)$\\
Also sind die Voraussetzungen von
\ref{EndlicheDarstellungDesSupremumsIntegrierbarerFunktionen} erfüllt, sodass
gilt:
\[\menge{\E\left(\inf_{i\in E}X_i\right)}{E\in\Pot_{ne}(I)}
 \overset{(\ast)}= -\menge{\E\left(\sup_{i\in E}Y_i\right)}{E\in\Pot_{ne}(I)}\]
ist nach \refclaim{EndlicheDarstellungDesSupremumsIntegrierbarerFunktionen}
{EndlicheDarstellungDesSupremumsIntegrierbarerFunktionen1} nach unten
beschränkt und
\begin{align*}
&\, \E\left( \inf_{i \in I} X_i\right)  \overset{(\ast)}{=}
 - \E\left( \sup_{i \in I} Y_i \right)
 \overset{\refclaim{EndlicheDarstellungDesSupremumsIntegrierbarerFunktionen}
{EndlicheDarstellungDesSupremumsIntegrierbarerFunktionen2}}=
- \sup\menge{\E\left(\sup_{i \in E}Y_i\right)}{E\in \Pot_{ne}(I)}\\
\overset{(\ast)}= &\,-\sup\menge{-\E\left(\inf_{i\in E}X_i\right)}{E\in\Pot_{ne}(I)}
 = \inf\menge{\E\left(\inf_{i\in E}X_i\right)}{E\in\Pot_{ne}(I)}\text.
\end{align*}
\end{proof}